% Included after the original top-order adjoint proof in the spectrum section.
\paragraph{Nine-moment, derivative-free form of the top-order adjoint.}
Using the notation of \S\ref{sec:cubic-moments}, put
\begin{align*}
 \mathscr P={}&-2\mathsf a n_h
 +2\langle\alpha_SS_2+\alpha_TT_2+2\mu\mathsf a n_2\rangle_\mu,\\
 \mathscr Q={}&2\mathsf a'n_h+2\left\langle
 \left[\frac{(1+2\mu)\mathsf b'}4q_S-\frac{\mathsf a'(3+\mu)}6\right]S_2\right.\\
 &\quad+\left[\frac{(1+2\mu)\mathsf b'}4q_T+\mathsf a'q_T\right]T_2\\
 &\quad\left.+\left[-(1+\mu)\mathsf b'(\mu-1/3)-2\mathsf a'(1+2\mu)\right]n_2
 \right\rangle_\mu .
\end{align*}
Then \eqref{eq:cubic-source-M} is exactly $\bar m=k^2(\mathscr P'+\mathscr Q)$.
Set $e=\bar t_\star$, $o=\bar t_o$, $w(t)=w_\ell(k(o-t))$ and
$\mathscr Z=\mathscr P-2\Upsilon$. Proposition~\ref{prop:adjoint} becomes
\begin{equation}\label{eq:adjoint-nine-moment}
\boxed{\begin{aligned}
 \mathcal A_\ell={}&14T_\ell+\int_0^o\mathcal W^0_\ell(s)\bar r_N(s)\dd s\\
 &+\frac32[w\mathscr Z]_e^o
 +\frac32\int_e^o\left[w\mathscr Q-w'\mathscr Z+\frac{2w}{t}\Upsilon\right]\dd t.
\end{aligned}}
\end{equation}
Here $w'$ differentiates $t$. All functions on the right require only the nine moments,
$n_h$, and known linear transfer profiles. Their relative-angle integrals are evaluated in
\eqref{eq:moment-N}--\eqref{eq:moment-V}; no cubic kernel or derivative of a second-order
response is needed. The boundary term includes the observer and must be kept.

To prove the identity, use
$\bar h_{S3}'=6(\bar\nu_3'-\bar\nu_3/t)+3\bar m/k^2$ and solve
$L_k\eta_3=-\bar r_N$ with regular homogeneous part $14\mathsf a$.
Recover $\bar\nu_3=\eta_3-\Upsilon$. Its contribution beyond the lapse Green integral is
$-3\int_e^ow(\Upsilon'-\Upsilon/t)\dd t$; the momentum contribution is
$\frac32\int_e^ow(\mathscr P'+\mathscr Q)\dd t$. Integrating the two derivatives by parts
gives \eqref{eq:adjoint-nine-moment}. The regular coefficient is unchanged since
$\Upsilon=O(t^4)$. The script \path{solution/check_shell_cubic_sources.py} checks the
source identities and the combined boundary coefficients on arbitrary jets.
This evaluates the relative-angle part of the top-order field response, not the complete
$R_\ell$: the lower-order propagation and measurement products must still be added.
